\documentclass[11pt]{article}
\usepackage{mathunal-base}
\mathunalfuente{serif}
\providecommand{\nota}[1]{\par\smallskip\noindent{\small\itshape\color{unalblue}#1}\par\smallskip}
\newlist{opciones}{enumerate}{1}
\setlist[opciones]{label=\alph*.,itemsep=1pt,topsep=2pt,leftmargin=1.8em,font=\normalfont}

\renewcommand{\examtitulo}{Métodos Numéricos 3002192 --- Primer Parcial}
\renewcommand{\examfecha}{Septiembre 4 de 2006}

\begin{document}

\examheader

\vspace{6pt}
\noindent Nombre \dotfill\ Carné \dotfill

\vspace{10pt}
\noindent\textbf{1.} Sea
\[
g(x)=-1+\frac{1}{4}(x+1)^2.
\]
\begin{opciones}
\item (5\%) Resuelva $x=g(x)$ y compruebe que los puntos fijos de $g(x)$ son $x=-1$ y $x=3$.
\item (5\%) Para el intervalo $I=[-2,0]$ pruebe que $g(x)\in I$ para todo $x\in I$. ¿Qué se concluye de este hecho?
\item (5\%) Para el mismo intervalo $I$, encuentre una constante $K$ tal que $|g'(x)|\leq K<1$. ¿Qué se concluye de este hecho?
\item (5\%) Realice dos iteraciones de punto fijo para $g(x)$ con punto inicial $x=-5$. Compruebe que se trata de una iteración convergente. ¿A qué punto fijo converge?
\item (5\%) Si el punto inicial es $x=3$, ¿será convergente la iteración de punto fijo? ¿A cuál punto fijo?
\item (5\%) Si el punto inicial es $x=-1.5$, ¿será convergente la iteración de punto fijo? ¿A cuál punto fijo?
\end{opciones}

\vspace{5mm}
\noindent\textbf{2.}
\begin{opciones}
\item (10\%) Sea $f(x)=x^{1/3}$. Enuncie el método iterativo de Newton Raphson para aproximar la solución de $f(x)=0$. Justifique la siguiente afirmación: La iteración de Newton Raphson para $f(x)=0$ converge únicamente si el punto inicial es la raíz $x=0$.
\item (10\%) Pruebe que la función $f(x)=x-\tan(x)$ tiene una raíz en $[-1.5,1]$. Encuentre dicha raíz. ¿Cuál es su multiplicidad? ¿Qué tipo de convergencia espera del método de Newton Raphson si lo aplica a este problema con punto inicial $x=-1$?
\end{opciones}

\vspace{5mm}
\noindent\textbf{3. (20\%)} Sea
\[
A=\begin{bmatrix} 1 & s \\ -s & 1 \end{bmatrix}.
\]
¿Bajo qué condiciones sobre $s$ convergen las iteraciones de Jacobi y Gauss Seidel?

\nota{Ayuda: $\begin{bmatrix} 1 & 0 \\ -s & 1 \end{bmatrix}^{-1}=\begin{bmatrix} 1 & 0 \\ s & 1 \end{bmatrix}$.}

\vspace{5mm}
\noindent\textbf{4.} Considere el sistema $Ax=b$, con $A$ no singular con diagonal no nula. Deseamos resolver este sistema por un método iterativo, ya sea Jacobi, Gauss Seidel o SOR. Suponga que $A=D-L-U$, donde $D$ es la matriz diagonal con la diagonal de $A$ como diagonal, $L$ es el negativo de la parte estrictamente triangular inferior de $A$ y $U$ es el negativo de la parte estrictamente triangular superior de $A$.
\begin{opciones}
\item (5\%) ¿Qué condición sobre la matriz $D^{-1}(L+U)$ garantiza la convergencia de uno de los tres métodos para cualquier vector inicial? ¿Cuál es el método que tiene que ser convergente en este caso?
\item (10\%) Establezca tres condiciones que garantizan la convergencia del método de Gauss Seidel para cualquier vector inicial.
\item (10\%) Suponga que $0<\omega<2$. Establezca dos condiciones que garantizan la convergencia del método SOR para cualquier vector inicial.
\item (5\%) Suponga que se sabe que la iteración de Gauss Seidel converge para cualquier vector inicial pero la iteración de Jacobi no converge para cualquier vector inicial. ¿Puede ser $A$ una matriz estrictamente diagonalmente dominante?
\end{opciones}

\vfill
\rule{\linewidth}{0.4pt}\\[1pt]
{\scriptsize\textit{Transcripción digital --- MathUNAL}\hfill\textcolor{unalblue}{\textbf{1}}}

\end{document}
